% === VIII. DISCUSSION AND LIMITATIONS ========================================
\section{Discussion and Limitations}
\label{sec:disc}

\subsection{What generalises, and what does not}

The reachable-friction bound~\eqref{eq:mureach} is kinematic: it involves no
tire parameter, no simulator-specific quantity and no property of the residual
model, so it applies to any implementation of the virtual-steady-state-data
strategy of~\cite{ontrack2025}, in simulation or on hardware. Its practical
content is a design rule: the rollout must be configured to demand at least the
friction one intends to identify, and where it cannot, $D$ should be reported
as unidentified rather than fitted. The diagnostic signature, coefficients
resting exactly on box constraints, is likewise generic and should be checked
and logged by any bounded tire-fitting routine, being otherwise
indistinguishable in the output from a successful fit. The two structural
properties of Section~\ref{subsec:contraction} generalise on the same grounds:
the $BCD$ degeneracy follows from the quasi-steady force recovery
\eqref{eq:ssforces}, and the absence of contraction from stage~4 alone. What
does not generalise is the value \num{0.20}, a property of this vehicle's
extrapolation error; what does is that $\beta=1$ is not defensible and
$\beta=0$ worse still. The insufficiency of per-coefficient box constraints is
also platform-independent, and the decaying surface adds a third instance to
the two we had: a stiffness that looks right while being the product of a
railed $B$ and a railed $D$, carried by a controller until the surface fell
below the reference's demand.

Three results are properties of this plant and must not be read as statements
about tires: the numerical tire values, since the plant is a PhysX
wheeled-vehicle model; the gap between configured wheel friction and realised
peak axle friction (\num{1.0} against \num{0.70}), which we report because an
identifier tuned to converge on the configuration file's number will converge
on the wrong one; and the \num{0.76} achieved/commanded steering ratio, where
the lesson (identify against the achieved actuator state) is general but the
magnitude is not. Sagmeister
\emph{et al.}~\cite{simfidelity2026} show sensitivity to model fidelity is
largest near the acceleration limits, where this study operates, and
R-CARLA~\cite{rcarla2025} exists because CARLA's native dynamics are a known
weak point for racing-grade work, so hardware replication is necessary before
any claim about identification accuracy on a real vehicle.

\subsection{Limitations}
\label{subsec:limits}

\begin{enumerate}
\item \textbf{One run per configuration and surface, so no metric in
Table~\ref{tab:comparison} carries an uncertainty.} The \SI{0.151}{\meter}
against \SI{0.224}{\meter} RMS tracking difference in particular has no
confidence interval and no evidence that it exceeds run-to-run variation, which
is not measured here; the order-of-magnitude force and friction margins are
large enough that such variation is unlikely to reverse them, but that is an
argument, not a measurement. The decaying-surface departure is a single event
and identifies a failure \emph{mechanism}, a grip ceiling pinned to a box
constraint while the plant falls below the reference's demand, not a failure
rate.
\item \textbf{One decay law, at one rate, on one vehicle and trajectory.} The
schedule is linear at \SI{0.1}{\percent} of nominal per second with a floor at
half, a factor \num{20} slower than the \SI{2}{\percent} per second
of~\cite{llampc2025}, for the reason Section~\ref{subsubsec:surface} gives: a
faster decay outruns the \SI{30}{\second} buffer this class of identifier is
built on, so the experiment would measure the sampling period rather than the
estimator. The runs therefore establish that the pipeline tracks a friction
change slow relative to its own cycle, with the scaling of
Section~\ref{subsec:runtime}. Step changes, which the scenario file also
carries, are not run, and the lag the decay exposes would appear on a step as
an overshoot of the whole cycle.
\item \textbf{The two surfaces were run on different controller builds.} The
decaying-surface pair carries a grip rate limiter, a binding-axle derate and the
fast-$\mu$ hook of Table~\ref{tab:scenarios}, none of which are in the
static-surface pair. Within each pair the two arms are identical apart from
Table~\ref{tab:scenarios}, so each comparison is clean; across pairs, only the
\oursarm{} lap-time and tracking-error trend is read.
\item \textbf{The friction estimator's excitation sweep is synthetic}
(Appendix~\ref{subsec:mu}), so the excitation dependence of its
\SI{11}{\percent} bias is not mapped against the simulator's per-wheel forces.
The published peak factor is exactly this estimate, so the bias passes through
undamped, and on the decaying surface it arrives with one cycle of lag on top
of it. Relatedly, the rear axle is not identified on this trajectory and the
front-axle $\mu$ is applied to both (Section~\ref{subsec:frictionfloor}); the
quasi-steady argument making $\mu_{f}=\mu_{r}$ holds on a homogeneous surface
and would not survive a split-$\mu$ one, which this sensor set cannot detect in
the first place.
\item \textbf{The two configurations differ in nine settings at once}
(Table~\ref{tab:scenarios}), so Section~\ref{subsec:summary} attributes nothing
to any one of them. Section~\ref{subsec:ablations} does attribute four, but by
offline replay on recorded buffers, so its ordering is not a closed-loop
ordering, and the integrator, the sweep bound and the regularisation are not
separated at all. The S4D figure there (\SI{21}{\percent} on force RMSE at
matched settings) has the same status: a measured ablation, not a reproduction
of the closed-loop S4 improvement of~\cite{visionsysid2026}. The
excitation-aware rollout, this paper's central contribution, is active in
\emph{both} arms, so the closed-loop comparison is not evidence for the
identifiability correction; the evidence for that is the sweep of
Section~\ref{subsec:sweep}.
\item \textbf{The physics-informed objective buys nothing measurable here.}
Once its units are reconciled with the data term and the loop is damped, it is
indistinguishable from a plain one-step MSE on this manoeuvre. It is retained
because it is the paper's stated formulation and costs nothing, not because it
is shown to help.
\item \textbf{Longitudinal dynamics are out of scope.} As in the baseline, only
lateral dynamics are identified. Combined slip and load transfer are neglected,
and the simulator's longitudinal wheel force channel is drivetrain torque rather
than a contact-patch force, so it is unusable as ground truth for a longitudinal
extension without further work.
\end{enumerate}
